\documentclass[fleqn]{article}
\oddsidemargin 0.0in
\textwidth 6.0in
\thispagestyle{empty}
\usepackage{import}
\usepackage{amsmath}
\usepackage{graphicx}
\usepackage{flexisym}
\usepackage{calligra}
\usepackage{amssymb}
\usepackage{bigints} 
\usepackage[english]{babel}
\usepackage[utf8x]{inputenc}
\usepackage{float}
\usepackage[colorinlistoftodos]{todonotes}


\DeclareMathAlphabet{\mathcalligra}{T1}{calligra}{m}{n}
\DeclareFontShape{T1}{calligra}{m}{n}{<->s*[2.2]callig15}{}
\newcommand{\scriptr}{\mathcalligra{r}\,}
\newcommand{\boldscriptr}{\pmb{\mathcalligra{r}}\,}

\definecolor{hwColor}{HTML}{1a0252}

\begin{document}

  \begin{titlepage}

    \newcommand{\HRule}{\rule{\linewidth}{0.5mm}}

    \center

    \begin{center}
      \includegraphics[height=11cm, width=11cm]{asu.png}
    \end{center}

    \vline

    \textsc{\LARGE Quantum Physics II}\\[1.5cm]

    \HRule \\[0.5cm]
    { \huge \bfseries Problem Set Eight}\\[0.4cm] 
    \HRule \\[1.0cm]

    \textbf{Behnam Amiri}

    \bigbreak

    \textbf{Prof: Onur Erten}

    \bigbreak

    \textbf{{\large \today}\\[2cm]}

    \vfill

  \end{titlepage}

  \begin{enumerate}
    \item A particle of spin $\dfrac{1}{2}$ is in a $p$ orbital $(\ell=1)$ of some potential.
    \begin{enumerate}
      \item What values of $J$, total angular momentum can it have?

        \textcolor{hwColor}{
          \\
          Knowing that the particle has $s=\dfrac{1}{2}, \ell=1$ and considering that $J=L+S$, by the 
          triangle rule gives $J=\left\{\dfrac{1}{2}, \dfrac{3}{2}\right\}$.
          \\
        }

      \item Use the method described in lecture (using the z-components and $\pm$ components of $J=L+S$
      to find the state $|j=1/2, j_z=1/2 \rangle$ in terms of the basis $|\ell=1, \ell_z \rangle|s=1/2, s_z \rangle$)
      Hint: First work out the $|j=3/2, j_z=3/2 \rangle$ then use orthogonality principle).

        \textcolor{hwColor}{
          \\
          $j=3/2$ and $j_z=3/2$ can be reached in one manner.
          \\
          \\
          $
            |j=3/2, j_z=3/2 \rangle=| \ell=1, \ell_z=1 \rangle ~ | s=1/2, s_z=1/2 \rangle
            \\
            \\
            \\
            J-|j=\dfrac{3}{2}, j_z=\dfrac{3}{2}|=\hbar \sqrt{J\bigg( j+1 \bigg)-j_z \bigg( j_z-1\bigg) } ~ | j, j_z-1 \rangle
            =\hbar ~ \sqrt{3} | j=3/2, j_z=1/2 \rangle
            \\
            \\
            \\
            \bigg( L_- + S_-\bigg) ~ | j=3/2, j_z=3/2 \rangle= \bigg( L_-+S_-\bigg) ~ | \ell=1, \ell_z=1 \rangle ~ | S=1/2, S_z=1/2 \rangle
            \\
            \\
            \\
            =\bigg( L_- |\ell=1, \ell_z=1 \rangle \bigg) ~ \bigg( S=1/2, S_z=1/2 \bigg)
            +|\ell=1, \ell_z=1 \rangle ~ \bigg( S_- | S=1/2, S_z=1/2 \rangle \bigg)
            \\
            \\
            \\
            =\hbar \sqrt{1+(1+1)-(1-1)} ~ |\ell=1, \ell_z=0 \rangle ~ | S=1/2, S_z=1/2 \rangle
            +| \ell=1 \rangle ~ | \ell_z=1 \rangle 
            \\
            \\
            \bigg( 
              \hbar \sqrt{1/2+(1/2+1)-1/2 (1/2-1)} ~ |S=1/2, S_z=1/2 \rangle  
            \bigg)
            \\
            \\
            \\
            \\
            \therefore ~~~ \sqrt{3} ~ \hbar ~ |\dfrac{3}{2}, \dfrac{1}{2} \rangle
            =\hbar ~ |1,1 \rangle | \dfrac{1}{2}, -\dfrac{1}{2} \rangle+\sqrt{2} \hbar |1,0 \rangle ~ |\dfrac{1}{2}, \dfrac{1}{2} \rangle
            \\
            \\
            \\
          $
          What we are looking for is $|\dfrac{1}{2}, \dfrac{1}{2} \rangle$. Posibilities with $j_z=\dfrac{1}{2}$ is that 
          $| \ell=1, \ell_z \rangle ~ |S=1/2, S_z$ basis are in $| \dfrac{3}{2}, \dfrac{1}{2} \rangle$ which are;
          \\
          \\
          $
            \begin{cases}
              |\ell=1, \ell_z=1 \rangle ~ |S=1/2, S_z=-1/2 \rangle
              \\
              \\
              |\ell=1, \ell_z=0 \rangle ~ |S=1/2, S_z=1/3 \rangle
            \end{cases}
            \\
            \\
          $
          $|j=1/2, j_z=1/2 \rangle$ is orthogonal and normalized to $|j=3/2, j_z=1/2 \rangle$ then we have:
          \\
          \\
          $
           \therefore ~~~ \boxed{
            |j=1/2, j_z=1/2 \rangle=\sqrt{\dfrac{3}{2}} ~ |\ell=1, \ell_z=1 \rangle ~ |S=1/2, S_z=-1/2 \rangle
            -\dfrac{1}{\sqrt{3}} ~ |\ell=1, \ell_z=0 \rangle ~ |S=1/2, S_z=1/2 \rangle
           } ~~~~ \checkmark
            \\
            \\
          $ 
        }

      % \bigg( \bigg)

      \item If the system is initially prepared in the state $|j=1/2, j_z=1/2 \rangle$, what is the 
      probability that measuring its value of $s_z$ would give $+1/2$?

        \textcolor{hwColor}{
          \\
          In the previous section we found for the state $|j=1/2, j_z=1/2 \rangle$ the basis. Therfore,
          the probability of measuring its value of $s_z$ gives $|-\dfrac{1}{\sqrt{3}}|^2=\dfrac{1}{3} ~~~~ \checkmark$.
          \\
          \\
        }

    \end{enumerate}

    \item \textbf{8.4, Zettili} Two identical particles of spin $\dfrac{1}{2}$ are enclosed in a one-dimensional 
    box potential of length $L$ with rigid walls at $x=0$ and $x=L$. Assuming that the two-particle system
    is in a triplet spin state, find the energy levels, the wave functions, and the degeneracies corresponding 
    to the three lowest states.

      % \textcolor{hwColor}{
      %   \\
      % }

  \end{enumerate}

\end{document}
